% GENERATED FILE. Edit canonical lessons, then run npm run generate:books.
% canonical-source-hash: fnv1a64:0b292db2e8d525eb
% canonical-lessons: ZH-C446-hutong, ZH-C446-cheng, ZH-W446-cheng-town, ZH-C446-gucheng, ZH-C446-chengqiang, ZH-C446-jincheng

\chapter{The Old Town}
\label{ch:zh-the-old-town}
\hlchaptermodality{446}

\begin{chapteropening}
\textbf{By the end of this chapter:} \emph{I can talk about old lanes, a walled town, an ancient town, its walls and moving to the city for work.}

Complete the last lesson of chapter 446: I can talk about old lanes, a walled town, an ancient town, its walls and moving to the city for work.
\end{chapteropening}

\section[hútòng]{hútòng — lane; alley (in an old city)}

\label{lesson:ZH-C446-hutong}



\begin{quote}
Before the new one: say the Mandarin for a seashell, then the Mandarin for to ebb.
\end{quote}

\subsection*{You'll want to know: hútòng}
\emph{hútòng} — \textquotedblleft{}lane; alley (in an old city)\textquotedblright{}.

\begin{quote}
\emph{wǒ zhù zài yì tiáo xiǎo hútòng lǐ} — I live in a little old lane.
\end{quote}

Its characters come later; for now it is a word for the ear.

\begin{cousinweb}
\emph{hútòng} names the narrow lanes of old northern cities. No dictionary we use records an older form, so its history is not settled.
\end{cousinweb}

\begin{grammarlens}[title={What you've built}]
One more word for getting around.
\end{grammarlens}

\subsection*{Your turn}
\begin{itemize}
\raggedright
  \item \emph{Say it:} \emph{hútòng}
  \item \emph{Say it:} \emph{hútòng}, once more
  \item \emph{Say it:} \emph{wǒ zhù zài yì tiáo xiǎo hútòng lǐ} — I live in a little old lane.
  \item \emph{Recall:} say the Mandarin for a seashell, then the Mandarin for to ebb, then say \emph{hútòng} again
\end{itemize}

\begin{tcolorbox}[breakable,colback=teal!4,colframe=teal!35!black,title={Before you move on}]
What does hútòng mean? (\textquotedblleft{}Lane; alley (in an old city)\textquotedblright{}.) Say it once more.
\end{tcolorbox}
\section[\texorpdfstring{\zh{城} (chéng)}{城 (chéng)}]{\zh{城} (chéng) — town; walled city}

\label{lesson:ZH-C446-cheng}



\begin{quote}
Before the new one: say the Mandarin for to ebb, then the Mandarin for a narrow old lane.
\end{quote}

\subsection*{You'll want to know: \zh{城}}
\textbf{\zh{城}} — \emph{chéng} — \textquotedblleft{}town; walled city\textquotedblright{}.

\begin{quote}
\textbf{\zh{城里有很多老房子}} \emph{chéng lǐ yǒu hěn duō lǎo fángzi} — There are lots of old houses in the town.
\end{quote}

The new character, in pieces: \zh{城} — nine strokes in two pieces: on the left the earth piece of \zh{地}, three strokes; on the right \zh{成} chéng, six strokes. The next lesson puts it in your hand.

\begin{cousinweb}
\textbf{\zh{城}} (\emph{chéng}) is inherited from Middle Chinese \emph{dzyeng} (Guangyun, in Baxter's transcription), which glosses it as the walls of a city and, after an older book, as what holds the people and their goods.
\end{cousinweb}

\begin{grammarlens}[title={What you've built}]
One more word for getting around.
\end{grammarlens}

\subsection*{Your turn}
\begin{itemize}
\raggedright
  \item \emph{Say it:} \emph{chéng}
  \item \emph{Say it:} \emph{chéng}, once more
  \item \emph{Say it:} \emph{chéng lǐ yǒu hěn duō lǎo fángzi} — There are lots of old houses in the town.
  \item \emph{Recall:} say the Mandarin for to ebb, then the Mandarin for a narrow old lane, then say \emph{chéng} again
\end{itemize}

\begin{tcolorbox}[breakable,colback=teal!4,colframe=teal!35!black,title={Before you move on}]
What does \zh{城} mean? (\textquotedblleft{}Town; walled city\textquotedblright{}.) Say it once more.
\end{tcolorbox}
\section[\texorpdfstring{\zh{城} (chéng)}{城 (chéng)}]{\zh{城} — the earth piece, then \zh{成}}

\label{lesson:ZH-W446-cheng-town}



\begin{quote}
The word first, then the character inside it.

\begin{itemize}
\raggedright
  \item \emph{Recall:} say \emph{chéng} — \textbf{R1}, one lesson back, before it decays
\end{itemize}

\begin{quote}
\zh{地} \zh{成}
\end{quote}

\emph{Read it:} each one aloud, with nothing said first — \emph{dì}, \emph{chéng}
\end{quote}

\subsection*{Script — the earth piece, then \zh{成}}
\hlblockfigure{\detokenize{figures/ZH-W446-cheng-town-filmstrip.pdf}}{How \zh{城} is written, stroke by stroke}

\begin{quote}
\zh{城}
\end{quote}

\textbf{Two pieces, left before right.} Count the pieces, not the strokes.

\textbf{The left piece, the earth piece, three strokes}, as on the left of \zh{地}: a short bar, a vertical, and a stroke that rises to the right.

\textbf{The right piece, \zh{成} chéng, six strokes}, as you wrote it: a bar; a long stroke falling to the left; a small bar that turns down and hooks; the long curved stroke with a hook; a short stroke falling to the left across it; and a dot at the top right.

3 + 6 = nine.

\subsection*{Your turn}
\begin{itemize}
\raggedright
  \item \emph{Write it:} \zh{城}, all nine, with the model in front of you
  \item \emph{Cover:} write it again from memory, and count the pieces, not the strokes
\end{itemize}

\begin{tcolorbox}[breakable,colback=teal!4,colframe=teal!35!black,title={Before you move on}]
How many strokes in \zh{城}? (\textbf{Nine}.) What is the left piece? (\textbf{The earth piece, as in \zh{地}}.) What is the right piece? (\textbf{\zh{成} chéng}.)

Source: \href{https://raw.githubusercontent.com/chanind/hanzi-writer-data/68d10a4b21150cae5e1ebbd223eed289cf32d90c/data/\%E5\%9F\%8E.json}{Hanzi Writer Data \zh{城}, pinned stroke paths and medians}.
\end{tcolorbox}
\section[\texorpdfstring{\zh{古城} (gǔchéng)}{古城 (gǔchéng)}]{\zh{古城} (gǔchéng) — ancient walled town; old town}

\label{lesson:ZH-C446-gucheng}



\begin{quote}
Before the new one: say the Mandarin for a narrow old lane, then the Mandarin for a walled town.
\end{quote}

\subsection*{You'll want to know: \zh{古城}}
\textbf{\zh{古城}} — \emph{gǔchéng} — \textquotedblleft{}ancient walled town; old town\textquotedblright{}.

\begin{quote}
\textbf{\zh{我们在古城里走了一个下午}} \emph{wǒmen zài gǔchéng lǐ zǒu le yí ge xiàwǔ} — We spent the afternoon walking round the old town.
\end{quote}

\begin{cousinweb}
\textbf{\zh{古城}} (\emph{gǔchéng}) is built from \zh{古} 'ancient' and \zh{城} 'a walled town', both inherited from Middle Chinese, \emph{kuX} and \emph{dzyeng} (Guangyun, in Baxter's transcription).
\end{cousinweb}

\begin{grammarlens}[title={What you've built}]
One more word for getting around.
\end{grammarlens}

\subsection*{Your turn}
\begin{itemize}
\raggedright
  \item \emph{Say it:} \emph{gǔchéng}
  \item \emph{Say it:} \emph{gǔchéng}, once more
  \item \emph{Say it:} \emph{wǒmen zài gǔchéng lǐ zǒu le yí ge xiàwǔ} — We spent the afternoon walking round the old town.
  \item \emph{Recall:} say the Mandarin for a narrow old lane, then the Mandarin for a walled town, then say \emph{gǔchéng} again
\end{itemize}

\begin{tcolorbox}[breakable,colback=teal!4,colframe=teal!35!black,title={Before you move on}]
What does \zh{古城} mean? (\textquotedblleft{}Ancient walled town; old town\textquotedblright{}.) Say it once more.
\end{tcolorbox}
\section[chéngqiáng]{chéngqiáng — city wall; the wall round a city}

\label{lesson:ZH-C446-chengqiang}



\begin{quote}
Before the new one: say the Mandarin for a walled town, then the Mandarin for an ancient walled town.
\end{quote}

\subsection*{You'll want to know: chéngqiáng}
\emph{chéngqiáng} — \textquotedblleft{}city wall; the wall round a city\textquotedblright{}.

\begin{quote}
\emph{qù chéngqiáng zěnme zǒu} — How do I get to the city wall?
\end{quote}

Its characters come later; for now it is a word for the ear.

\begin{cousinweb}
\emph{chéngqiáng} is built from \zh{城} 'a walled town' and \emph{qiáng} 'a wall', both inherited from Middle Chinese, \emph{dzyeng} and \emph{dzjang} (Guangyun, in Baxter's transcription).
\end{cousinweb}

\begin{grammarlens}[title={What you've built}]
One more word for getting around.
\end{grammarlens}

\subsection*{Your turn}
\begin{itemize}
\raggedright
  \item \emph{Say it:} \emph{chéngqiáng}
  \item \emph{Say it:} \emph{chéngqiáng}, once more
  \item \emph{Say it:} \emph{qù chéngqiáng zěnme zǒu} — How do I get to the city wall?
  \item \emph{Recall:} say the Mandarin for a walled town, then the Mandarin for an ancient walled town, then say \emph{chéngqiáng} again
\end{itemize}

\begin{tcolorbox}[breakable,colback=teal!4,colframe=teal!35!black,title={Before you move on}]
What does chéngqiáng mean? (\textquotedblleft{}City wall; the wall round a city\textquotedblright{}.) Say it once more.
\end{tcolorbox}
\section[\texorpdfstring{\zh{进城} (jìnchéng)}{进城 (jìnchéng)}]{\zh{进城} (jìnchéng) — to go to town; to move to the city}

\label{lesson:ZH-C446-jincheng}



\begin{quote}
Before the new one: say the Mandarin for an ancient walled town, then the Mandarin for the wall round a city.
\end{quote}

\subsection*{You'll want to know: \zh{进城}}
\textbf{\zh{进城}} — \emph{jìnchéng} — \textquotedblleft{}to go to town; to move to the city\textquotedblright{}.

\begin{quote}
\textbf{\zh{村子里很多人进城打工}} \emph{cūnzi lǐ hěn duō rén jìnchéng dǎgōng} — A lot of people from the village move to the city for work.
\end{quote}

\begin{cousinweb}
\textbf{\zh{进城}} (\emph{jìnchéng}) is built from \zh{进} 'to enter' and \zh{城} 'a walled town', both inherited from Middle Chinese, \emph{tsinH} and \emph{dzyeng} (Guangyun, in Baxter's transcription), which glosses the first 'to go forward'.
\end{cousinweb}

\begin{grammarlens}[title={What you've built}]
One more word for getting around.
\end{grammarlens}

\subsection*{Your turn}
\begin{itemize}
\raggedright
  \item \emph{Say it:} \emph{jìnchéng}
  \item \emph{Say it:} \emph{jìnchéng}, once more
  \item \emph{Say it:} \emph{cūnzi lǐ hěn duō rén jìnchéng dǎgōng} — A lot of people from the village move to the city for work.
  \item \emph{Recall:} say the Mandarin for an ancient walled town, then the Mandarin for the wall round a city, then say \emph{jìnchéng} again
\end{itemize}

\begin{tcolorbox}[breakable,colback=teal!4,colframe=teal!35!black,title={Before you move on}]
What does \zh{进城} mean? (\textquotedblleft{}To go to town; to move to the city\textquotedblright{}.) Say it once more.
\end{tcolorbox}
